\section{Méthodologie}\label{sec-methode}

\subsection{Jeu de données}\label{sec-donnees}

\paragraph{Source.}
Nous utilisons BMD-45~\cite{sharma2026bmd45}, un jeu de données publié à
CVPR 2026 (Findings) et diffusé sous licence CC~BY~4.0. Il comprend
45\,000 images et 480\,000 boîtes englobantes annotées, issues de plus de
3\,600 caméras CCTV opérationnelles du réseau de vidéosurveillance de
Bengaluru (Inde), avec 14 catégories fines de véhicules. Ce jeu présente
trois caractéristiques qui motivent son choix : des caméras fixes de
vidéoprotection (et non des caméras embarquées), un trafic urbain dense,
et une forte proportion de deux-roues, trois traits structurels partagés
avec le trafic urbain béninois (cf.~\cref{sec-discussion} pour les limites
de ce rapprochement).

\paragraph{Extraction du sous-ensemble.}
Pour tenir dans le budget de calcul disponible, un sous-ensemble est
extrait en streaming depuis Hugging Face avec des graines fixées :
2\,400 images d'entraînement (graine 42) et 600 images de validation
(graine 43). BMD-45 ne fournissant pas de jeu de test, celui-ci est dérivé
de la validation par une partition déterministe moitié-moitié (graine 42),
soit au final 2\,400 images d'entraînement, 300 de validation et 300 de
test. Le jeu de test n'intervient à aucun moment dans l'entraînement ni
dans la sélection du modèle.

\paragraph{Remappage des classes.}
Le modèle de référence étant pré-entraîné sur COCO~\cite{lin2014coco},
la comparaison n'est équitable que si les deux modèles partagent le même
espace sémantique. Les 14 classes de BMD-45 sont donc ramenées aux
4 classes de véhicules communes avec COCO selon le
\cref{tab-remappage}. Les trois catégories sans équivalent COCO
(\emph{Three-wheeler}, \emph{Bicycle}, \emph{Other}) sont ignorées :
les conserver reviendrait à demander au modèle pré-entraîné de détecter
des classes qu'il ne connaît pas, ce qui fausserait la comparaison en sa
défaveur.

\begin{table}[tb]
  \centering
  \caption{Remappage des 14 classes BMD-45 vers les 4 classes alignées
  sur COCO. L'identifiant COCO est celui utilisé lors de la réindexation
  du protocole d'évaluation (\cref{sec-protocole}).}
  \label{tab-remappage}
  \begin{tabular}{lll}
    \toprule
    Classes BMD-45 & Classe retenue & Identifiant COCO \\
    \midrule
    Hatchback, Sedan, SUV, MUV, Van & voiture & 2 (\emph{car}) \\
    Two-wheeler                     & moto    & 3 (\emph{motorcycle}) \\
    Bus, Mini-bus, Tempo-traveller  & bus     & 5 (\emph{bus}) \\
    Truck, LCV                      & camion  & 7 (\emph{truck}) \\
    \midrule
    Three-wheeler, Bicycle, Other   & \multicolumn{2}{l}{ignorées (aucun équivalent COCO)} \\
    \bottomrule
  \end{tabular}
\end{table}

\paragraph{Volumétrie.}
Le \cref{tab-volumetrie} donne les effectifs par split et par classe
après remappage. Deux faits saillants : le sous-ensemble compte en
moyenne $8{,}5$ véhicules annotés par image, signature d'un trafic dense
propice aux occlusions ; et les motos représentent $56\,\%$ des
25\,688 instances, une distribution qui rapproche structurellement ce jeu
du trafic béninois, dominé par les taxis-motos (zémidjans).

\begin{table}[tb]
  \centering
  \caption{Effectifs du sous-ensemble par split et par classe après
  remappage vers les 4 classes alignées sur COCO.}
  \label{tab-volumetrie}
  \begin{tabular}{lrrrrrr}
    \toprule
    Split & Images & voiture & moto & bus & camion & Total instances \\
    \midrule
    train & 2\,400 & 5\,744 & 11\,641 & 1\,298 & 1\,873 & 20\,556 \\
    val   & 300    & 752    & 1\,476  & 138    & 215    & 2\,581  \\
    test  & 300    & 742    & 1\,380  & 189    & 240    & 2\,551  \\
    \midrule
    Total & 3\,000 & 7\,238 & 14\,497 & 1\,625 & 2\,328 & 25\,688 \\
    \bottomrule
  \end{tabular}
\end{table}

\subsection{Modèle et entraînement}\label{sec-entrainement}

\paragraph{Modèle.}
Nous partons de YOLOv8n~\cite{yolov8_ultralytics}, la déclinaison la plus
compacte ($3{,}2$ millions de paramètres) de la famille
YOLOv8, héritière de l'architecture à une passe introduite par
YOLO~\cite{redmon2016yolo}. Le choix du modèle \emph{nano} est un
compromis assumé entre coût d'entraînement et représentativité : c'est
aussi la taille la plus plausible pour une inférence temps réel sur du
matériel modeste. Le fine-tuning démarre des poids pré-entraînés sur
COCO~\cite{lin2014coco}, qui servent également de modèle de référence
dans la comparaison.

\paragraph{Hyperparamètres.}
L'entraînement utilise Ultralytics 8.4.114 et PyTorch 2.13, sur un GPU
NVIDIA T4 (Google Colab) : 50 epochs, images redimensionnées à 640 pixels,
batch de 16, arrêt anticipé avec patience de 15 epochs, graine aléatoire 42.
Le \cref{tab-hyperparams} récapitule les augmentations de données.
Deux d'entre elles sont directement rattachées à la problématique :
\emph{mosaic} (assemblage de quatre images en une) et \emph{scale}
(zoom aléatoire $\pm 70\,\%$) produisent mécaniquement des objets plus
petits que ceux du jeu d'origine, ce qui expose le modèle à la
configuration visée, celle des véhicules éloignés et de petite taille. À
l'inverse, le retournement vertical et la rotation sont désactivés
(\texttt{flipud}${}=0$, \texttt{degrees}${}=0$) : les caméras de
vidéoprotection sont fixes et une scène routière possède un haut et un
bas ; ces transformations créeraient des configurations jamais
rencontrées en exploitation.

\begin{table}[tb]
  \centering
  \caption{Hyperparamètres d'augmentation de données. Les valeurs non
  listées sont celles par défaut d'Ultralytics 8.4.114.}
  \label{tab-hyperparams}
  \begin{tabular}{lll}
    \toprule
    Paramètre & Valeur & Justification \\
    \midrule
    \texttt{mosaic}    & $1{,}0$   & produit des objets plus petits (problématique) \\
    \texttt{scale}     & $0{,}7$   & zoom aléatoire, mêmes effets \\
    \texttt{fliplr}    & $0{,}5$   & symétrie gauche-droite plausible \\
    \texttt{translate} & $0{,}1$   & translations légères \\
    \texttt{hsv\_h}    & $0{,}015$ & variations de teinte \\
    \texttt{hsv\_s}    & $0{,}7$   & variations de saturation \\
    \texttt{hsv\_v}    & $0{,}4$   & variations de luminosité (heures, météo) \\
    \texttt{flipud}    & $0{,}0$   & désactivé : une scène routière a un haut et un bas \\
    \texttt{degrees}   & $0{,}0$   & désactivé : caméras fixes \\
    \bottomrule
  \end{tabular}
\end{table}

\subsection{Protocole d'évaluation}\label{sec-protocole}

Les deux modèles (YOLOv8n pré-entraîné sur COCO et YOLOv8n fine-tuné)
sont évalués sur exactement les mêmes 300 images de test et les mêmes
2\,551 boîtes de référence. Trois précautions garantissent l'équité et la
pertinence de la comparaison.

\paragraph{1. Alignement et réindexation des classes.}
Le modèle pré-entraîné expose les 80 classes COCO, le modèle fine-tuné
en expose 4. Les labels du jeu de test sont donc réindexés à la volée
vers l'espace de classes du modèle évalué : pour le modèle COCO,
$0 \rightarrow 2$, $1 \rightarrow 3$, $2 \rightarrow 5$,
$3 \rightarrow 7$ ; pour le modèle fine-tuné, les indices $0$ à $3$ sont
conservés. Le modèle COCO est de plus restreint aux classes
$\{2, 3, 5, 7\}$ lors de l'inférence, afin qu'il ne soit pas pénalisé
par des détections hors périmètre (piétons, feux de signalisation, etc.)
qui seraient comptées comme faux positifs.

\paragraph{2. Rappel ventilé par taille d'objet.}
La problématique portant sur les véhicules petits et éloignés, le rappel
est ventilé selon la convention COCO : \emph{petit} si l'aire de la boîte
est inférieure à $32 \times 32$ pixels, \emph{moyen} si elle est
inférieure à $96 \times 96$ pixels, \emph{grand} au-delà. Point crucial :
les aires sont ramenées à la résolution d'entrée du réseau (640 pixels)
avant classement. Sans cette normalisation, sur des images de plusieurs
milliers de pixels de large, la quasi-totalité des véhicules tombe dans
la catégorie \og grand \fg{} et la ventilation ne mesure plus rien.
C'est précisément ce qui s'est produit lors de notre premier essai, et ce
qui a conduit à cette correction du protocole. Le calcul procède par
appariement glouton : les prédictions (seuil de confiance $0{,}25$) sont
appariées aux boîtes de référence par IoU décroissant
(\cref{eq-iou}), avec un seuil $\mathrm{IoU} \geq 0{,}5$ ; une
prédiction ne peut valider qu'une seule boîte de référence ; la mesure
est agnostique à la classe.

\paragraph{3. Jeu de test isolé.}
Le jeu de test, dérivé de façon déterministe
(\cref{sec-donnees}), n'a jamais été vu à l'entraînement ni utilisé pour
la sélection du modèle ou des hyperparamètres.

\paragraph{Métriques.}
Nous rapportons la précision, le rappel et le F1-score
(\cref{eq-precision-rappel,eq-f1}), la mAP à IoU fixé à $0{,}5$ et
moyennée de $0{,}5$ à $0{,}95$ (\cref{eq-map}), le débit d'inférence en
images par seconde (FPS, mesuré sur GPU T4, une inférence d'échauffement
exclue de la mesure), et le rappel ventilé petit / moyen / grand décrit
ci-dessus. Les métriques globales sont calculées par le validateur
d'Ultralytics avec un seuil de confiance de $0{,}001$ et un seuil IoU de
$0{,}6$ : le seuil de confiance très bas est nécessaire pour tracer la
courbe précision-rappel complète dont la mAP est l'aire sous la courbe.

\begin{equation}
  \text{Précision} = \frac{VP}{VP + FP}, \qquad
  \text{Rappel} = \frac{VP}{VP + FN},
  \label{eq-precision-rappel}
\end{equation}
où $VP$, $FP$ et $FN$ désignent respectivement les vrais positifs, faux
positifs et faux négatifs ;
\begin{equation}
  F_1 = 2 \cdot \frac{\text{Précision} \times \text{Rappel}}
                     {\text{Précision} + \text{Rappel}} ;
  \label{eq-f1}
\end{equation}
\begin{equation}
  \mathrm{IoU}(A, B) = \frac{|A \cap B|}{|A \cup B|}
  \label{eq-iou}
\end{equation}
pour deux boîtes $A$ et $B$ ; et
\begin{equation}
  \mathrm{mAP} = \frac{1}{N} \sum_{i=1}^{N} \mathrm{AP}_i,
  \label{eq-map}
\end{equation}
où $\mathrm{AP}_i$ est l'aire sous la courbe précision-rappel de la
classe $i$ et $N$ le nombre de classes.
